\section{Conclusion}
\label{sec:conclusion}

SBOM specifications can accept structurally valid artifacts while leaving important generator choices unclear to consumers. In the evaluated CycloneDX outputs, the mandatory empty-dependency-entry rule was not schema-enforced and was violated by 13 applicable cdxgen outputs and 3 Syft outputs. No evaluated CycloneDX output populated document-level supplier metadata or declared composition-based coverage. The selected SPDX rules produced no Unmet result, providing a useful control: emitted constructs satisfied those checks, while N/A constructs were simply not exercised. Across both formats, profiles and rich data models represent provenance, relationships, vulnerabilities, build context, and coverage, but they do not provide a uniform cross-field vocabulary for explaining absent data.

The paper contributes a 2,286-row requirement corpus and clause-to-source-to-output evidence chains that support the cumulative suite in Section~\ref{sec:root-causes}. Level~1 checks identity and provenance, Level~2 checks evidence, filtering, coverage, and absence, and Level~3 evaluates vulnerability/VEX and build context against task needs. The four frozen snapshots and generated SBOMs trace these mechanisms rather than estimate ecosystem-wide prevalence. Implementing the suite and measuring consumer outcomes remain future work.
